% ============ §0 讲义档案 ============
\doctitle{俄语　全课程精讲}
{从 33 个字母到「三时态 ＋ 数词」：一册能脱离课件自学的入门讲义}
{依据：УРОК 1--8（8 份）＋ 动词现在／过去／将来时（3 份）＋ 数词（1 份），共 12 份课件、344 页 \quad 2026-10-23}

% ============ 目录（手排，不用 \tableofcontents，因为节头是自定义宏） ============
\subhdr{本册目录}

\begin{minipage}{\linewidth}
\small
\begin{tabular}{@{}>{\raggedright\arraybackslash}p{0.08\linewidth}>{\raggedright\arraybackslash}p{0.30\linewidth}>{\raggedright\arraybackslash}p{0.56\linewidth}@{}}
\S0 & 讲义档案 & 12 份课件的坐标、三种标记、读法、页数交代 \\
\S1 & 语音与字母 & 33 个字母分三类；元音五对；重音与 о 弱化；硬软辅音；清浊与两条清化；三个调型；五个基本问答 \\
\S2 & 名词 & кто／что；三个性与 -ь 的 38 词；复数三表合一；24 组特殊复数；只有单／复数 \\
\S3 & 人称与物主代词 & 八个主语；ты／Вы；四列式物主代词；先判性再选形；Меня́ зову́т 与问候 \\
\S4 & 形容词 & 位置与 како́й 四形式；五族词尾总表；拼写规则七字母；13 个颜色词；三步与四个坑 \\
\S5 & 指示代词与 нра́виться & э́тот 四形式；са́мый；нра́виться ＋ 第三格；что／потому́ что；ИК-5；课文 ПАРК СИНХА́Й \\
\S6 & 动词 & 两个变位组；-ся 动词；-л 型过去时；бу́ду ＋ 不定式；三时态对照；по-ру́сски 与副词 \\
\S7 & 数词 & 1--1000；读长数字；价格（юа́нь 三形）；时间（час 三形）；序数词；星期与月份 \\
\S8 & 复习卡 & 四张速查表：名词／词尾／动词／数词 \\
\S9 & 易错点 & 十二条，每条都能指回某个词的某一步 \\
\S10 & 缺口清单 & 12 条：课件跳了哪一步，本讲义补在哪 \\
\S11 & 三句话总结 & 判性数、一条规律多个位置、能回指页码 \\
\S12 & 页码对照表 & 12 份课件 → 本讲义哪一节 \\
\end{tabular}
\end{minipage}

\vspace{2pt}

{\small\color{inkgray}读法建议：第一次顺着 \S1 到 \S7 读；考前只读 \S8 与 \S9；做练习卡住时先查 \S8 的表，再回对应节的「关键一步」。}


\hdr{§0}{这份讲义讲什么、按什么顺序、怎么读}

\callout{\textbf{讲义档案}\\[3pt]
\textbf{讲什么}：老师 12 份课件里的全部正课内容——字母与语音（У1--У3）、名词的性与数（У4--У5）、人称与物主代词（У6）、形容词与颜色（У7）、指示代词与 \textbf{нра́виться}（У8）、动词三个时态、数词 1--1000 与「几点／多少钱」。\\
\textbf{不重复什么}：课件里成串的朗读操练页（同一个句型列十几行）不逐页重打，只在需要时举一行作代表；具体删了哪些，\S0.4 有一条一条的交代。\\
\textbf{补什么}：课件默认你会补的那些「为什么」——为什么 \textbf{я} 让前面的辅音变软、为什么非重音的 \textbf{о} 读成 [а]、为什么 \textbf{друг} 读 [друк]、为什么 \textbf{-ь} 既可能是阳性也可能是阴性、为什么 \textbf{г/к/х、ж/ш/ч/щ} 后面只能写 \textbf{и}。每条都在 \S10 点名页码。}

\subhdr{0.1　页码坐标怎么认}

全篇只有一个规矩：\textbf{每个结论都能回指某一页}。写法是「课件简称 ＋ 页数」，页数用 PDF 页码（不是胶片右下角那个编号，两者偶尔差一页，本讲义一律以 PDF 页码为准）。

\begin{minipage}{\linewidth}
\tabtop
\begin{tabular}{@{}>{\raggedright\arraybackslash}p{0.20\linewidth}>{\raggedright\arraybackslash}p{0.46\linewidth}>{\raggedright\arraybackslash}p{0.26\linewidth}@{}}
\thead{简称} & \thead{对应课件} & \thead{页数} \\
\end{tabular}
\tabmid
\begin{tabular}{@{}>{\raggedright\arraybackslash}p{0.20\linewidth}>{\raggedright\arraybackslash}p{0.46\linewidth}>{\raggedright\arraybackslash}p{0.26\linewidth}@{}}
У1 & УРОК 1. Алфавит & 41 \\
У2 & УРОК 2. Твёрдые и мягкие согласные звуки. Оглушение. Озвончение. Диалоги & 37 \\
У3 & УРОК 3. Глухие и звонкие согласные звуки. Ритмика. ИК-3 & 23 \\
У4 & УРОК 4. Имя существительное. Род & 19 \\
У5 & УРОК 5. Множественное число существительных & 14 \\
У6 & УРОК 6. Личные и притяжательные местоимения & 27 \\
У7 & УРОК 7. Имя прилагательное. 形容词 & 52 \\
У8 & УРОК 8. Указательные местоимения. Нравиться & 30 \\
动现 & ГЛАГОЛ. НАСТОЯЩЕЕ ВРЕМЯ & 37 \\
动过 & Глагол. Прошедшее время & 24 \\
动将 & ГЛАГОЛ. БУДУЩЕЕ ВРЕМЯ & 15 \\
数词 & 1 (оди́н) 2 (два) … 10 (де́сять) & 25 \\
\end{tabular}
\tabbot
\end{minipage}

\textbf{这张表怎么用}：读到「\textbf{У4 p13} 左栏」就翻 УРОК 4 那份 PDF 的第 13 页；读到「动现 p5」就翻动词现在时那 37 页里的第 5 页。12 份课件合计 344 页，本讲义每一节末尾都标了它对应哪几页。

\subhdr{0.2　三种标记}

\begin{itemize}
\item \textbf{重音撇}：вода́、хорошо́、молоко́——撇在哪个元音上，那个音节就读得最重、最长。俄语一个词通常只有一个重音；重音位置换了，词可能就是另一个词（如 \textbf{мука́} 面粉 ／ \textbf{му́ка} 痛苦），所以本讲义把它当拼写的一部分照抄。
\item \textbf{方括号}：[друк]、[шы]、[фк]——方括号里是\textbf{实际读音}，不是写法。写法 \textbf{друг}、读音 [друк]，就是 \S1.5 那条清化规律在起作用。
\item \textbf{星号}：*ру́скый、*э́тот пальто́——星号是\textbf{错写法}，专门用来标「很容易写成这样，但不对」。
\end{itemize}

\aside{为什么重音要跟着抄：俄语的弱化、清化、词尾变化都看重音落在哪。少了那一撇，后面的语音规则（\S1.3、\S1.5）就无从判断。}

\subhdr{0.3　本讲义的三条读法}

\begin{enumerate}
\item \textbf{第一次读}：§1 到 §7 顺着读。每节的「卡」是表格，「讲」是表格上下的文字——文字不是装饰，它解释这张表怎么来的、每一格为什么是那个值。只看表不看字，等于把课件又过了一遍。
\item \textbf{考前}：只读 \S8 复习卡 ＋ \S9 易错点。前者是全部词尾的一张速查，后者每条都指回某个具体词的某一步。
\item \textbf{写作业卡住时}：先查 \S8 的表，再查对应的节；每节末尾的练习都给了「关键一步」，卡在哪一步就看那一步为什么这么做。
\end{enumerate}

\subhdr{0.4　页数交代：为什么本讲义比课件短}

源课件 344 页，其中 \textbf{116 页没有文字层}（整页是图片：字母卡片、朗读操练、人物照片），另有大量页面是同一句型的重复操练。本讲义不逐页重打，做了三类压缩，这里逐条交代：

\begin{itemize}
\item \textbf{朗读操练页合并}：У1 p19--p25 的「Кто э́то? ／ Что э́то?」左右对照，压成 \S1.7 的一组问答加一张对照；У2 p10--p12 的三列音节练习，压成 \S1.4 的一张软硬对照表。
\item \textbf{图片页只取骨架}：字母卡片页（У1 p9 的字母表、У3 p3 的清浊词表）按图片实际内容转录，转录件在正文里注明「据图页转录」；照片、听录音页不转录，只在需要时点名页码。
\item \textbf{重复的例词只留一行}：同一个词在课件里出现四五次（如 \textbf{стол} 在 У1 出现 3 次、У7 出现 2 次），本讲义只在第一次出现处展开。
\end{itemize}

\textbf{没有删的是}：每个语法点的规则本身、每个不规则词形（\textbf{-ь} 的 38 个词、18 组特殊复数、13 个颜色词、两个变位的全部词尾、序数词 1--10）、课件的例题与练习。这三类在本讲义里只有增加，没有减少。
